\begin{abstract}
% 中文: LLM 在不同代码原语上表现出异质性，但 accuracy 无法解释内部差异来源。我们区分
% "信息已存在"与"信息已可用"，构建 dual diagnostic framework 追踪代码推理的内部生命周期。
% 发现 brewing 过程和四种 resolution outcome，不同 code primitive 诱导不同 fingerprint，
% brewing scaffold 跨模型稳定而 outcome 随能力变化。

LLMs show pronounced heterogeneity on code reasoning, yet task accuracy alone cannot explain why---similar outputs can arise from different internal failures. We distinguish \textit{information availability} (whether the answer is present in a hidden state) from \textit{information readiness} (whether the model can stably decode it), and introduce a dual diagnostic framework combining layer-wise linear probing with Context-Stripped Decoding across six code task families and five model families (0.5B--14B). We find a consistent internal lifecycle: models first pass through a \textbf{brewing} stage where the answer is linearly recoverable before it becomes self-decodable, then enter one of four causally validated \textbf{resolution outcomes}---Resolved, Overprocessed, Misresolved, or Unresolved. Most primitives are predominantly Resolved, but Computing stands out below 34\% and is the only task where all four outcomes coexist substantially. Controlled difficulty sweeps reveal that other primitives harbor latent failure bottlenecks that surface only under increased complexity. Across architectures and scales, the brewing scaffold remains stable while resolution success changes, suggesting that mechanism and capability are partly separable.
\end{abstract}
